\section{Work}
%%%%%%%%%%%%%%%%%%%%%%%%%%%%%%%%%%%%%%%%%%%%%%%%%%%%%%%%%%%%%%%%%
\subsection{The System Model}
%%%%%%%%%%%%%%%%%%%%%%%%%%%%%%%%%%%%%%%%%%%%%%%%%%%%%%%%%%%%%%%%%
Let $x\in\mathbb{R}^{n\times 1}$, $y\in\mathbb{R}^{p\times 1}$ and $w\in\mathbb{R}^{r\times 1}$ be the state, measurement, exogenous vectors, respectively, the SIMO/MIMO LTI system addressed in this paper is stated as, 
\begin{gather} 
    \dot{x}=A x+B_w w,\,y=C x+D_w w
\end{gather}
where $A\in\mathbb{R}^{n\times n}$, $B_w\in\mathbb{R}^{n\times r}$, $C\in\mathbb{R}^{p\times n}$ and $D_w\in\mathbb{R}^{p\times r}$. The following additional rank condition
\begin{gather}
    \mathrm{rank}(C)=r<p
\end{gather}
arises in sparse sensor applications\cite{shoukry2015event}, topologies used in Multi Agen Systems(MAS)\cite{olfati2007consensus} and distributed networks \cite{yang2022sensor}, and is also called Strictly Output Redundant(SOR) system \cite{yang2025output}. Here, the system output is overdetermined, therefore, 
\begin{gather}  
    \mathrm{dim}\,\mathcal{N}(C^T)=p-r\geq 1
\end{gather}  
or using orthogonality $C^TC^\perp=0$,
\begin{gather}  
    \mathrm{dim}\,\mathcal{N}(C^\perp)=p-r\geq 1
\end{gather}  
which ensures nontrivial orthogonal component in the output space. 
%%%%%%%%%%%%%%%%%%%%%%%%%%%%%%%%%%%%%%%%%%%%%%%%%%%%%%%%%%%%%%%%%
\subsection{The Observer}
%%%%%%%%%%%%%%%%%%%%%%%%%%%%%%%%%%%%%%%%%%%%%%%%%%%%%%%%%%%%%%%%%
The Luenberger Observer is defined as,    
\begin{gather} 
    \dot{\hat{x}}=A \hat{x}+L(y-\hat{y}),\,\hat{y}=C \hat{x}
\end{gather}
where $L\in\mathbb{R}^{n\times p}$ is the observer gain. Assuming $(A,C)$-pair observable, the error dynamics are obtained for the error $e\triangleq x-\hat{x}$ as,
\begin{gather} 
    \dot{e}=(A-LC)e+(B_w-L D_w)w
\end{gather}
The $\mathcal{H}_\infty$ optimal observer for objective function $z=C_ze$ is designed using the following LMI problem\cite{duan2013lmis},
\begin{equation*} \label{eq:p_lmi1}
\begin{split} 
    \min_{Y,P}{(\gamma)}\quad \mathrm{s.t.}\\
    \begin{bmatrix}
        (PA-YC)+\star& PB_w-YD_w& C_z^T\\
        \star& -\gamma I& 0\\
        \star& 0& -\gamma I
    \end{bmatrix}\prec 0\\
    P\succ 0
\end{split} 
\end{equation*}
where the observer gain is recovered with $L=P^{-1}Y$. 

%%%%%%%%%%%%%%%%%%%%%%%%%%%%%%%%%%%%%%%%%%%%%%%%%%%%%%%%%%%%%%%%%
\subsection{The Projection}
%%%%%%%%%%%%%%%%%%%%%%%%%%%%%%%%%%%%%%%%%%%%%%%%%%%%%%%%%%%%%%%%%
The projection operator defined using Gauge-Invariance principle is "proposed" as follows,
\begin{gather}
    \Pi\triangleq I-C(C^TC)^{-1}C^T
\end{gather}
which satisfies the following properties,
\begin{gather}
\begin{split}
    &\Pi C= (I-C(C^TC)^{-1}C^T)C=0\\
    &\Pi C^\perp=(I-C(C^TC)^{-1}C^T)C^\perp=C^\perp,\,C^TC^\perp=0\\
    &\Pi^2=\Pi,\,\Pi\in\mathbb{R}^{p\times p}
\end{split}
\end{gather}
Using the projection operator on the measured output yields,
\begin{gather}
    \Pi y=\Pi(C x+D_w w)=\Pi D_w w
\end{gather}
The projection eliminates the state dependent component of the output, isolating the exogenous component establishing state invariance. This is an algebraic measurement of the disturbance reflected to the output.The exogenous input is recovered under the following condition,
\begin{gather}
    \mathrm{rank}(\Pi D_w)=r
\end{gather}
or equivalently,
\begin{gather}
\ker(\Pi D_w) = {0}.
\end{gather}
with the least square estimator,
\begin{gather}
    \hat{\omega}=(\Pi D_w)^\dagger(\Pi y)
\end{gather}
where $(.)^\dagger$ is the Moore-Penrose pseudoinverse. The exogenous signal recovery fails if 
\begin{gather}
D_w w \in \mathrm{Im}(C),
\end{gather}
%%%%%%%%%%%%%%%%%%%%%%%%%%%%%%%%%%%%%%%%%%%%%%%%%%%%%%%%%%%%%%%%%
\subsection{The GI-Luenberger Observer}
%%%%%%%%%%%%%%%%%%%%%%%%%%%%%%%%%%%%%%%%%%%%%%%%%%%%%%%%%%%%%%%%%
The isolated disturbance is added to the classical observer as follows,
\begin{gather} 
    \dot{\hat{x}}=A \hat{x}+L(y-\hat{y})+B_w\hat{w},\,\hat{y}=C \hat{x}+D_w\hat{w}\\
    \hat{w}=(\Pi D_w)^\dagger(\Pi y)
\end{gather}
where the observer estimates the states and the projection the disturbance. 



%%%%%%%%%%%%%%%%%%%%%%%%%%%%%%%%%%%%%%%%%%%%%%%%%%%%%%%%%%%%%%%%%
\section{Projection Filtering}
%%%%%%%%%%%%%%%%%%%%%%%%%%%%%%%%%%%%%%%%%%%%%%%%%%%%%%%%%%%%%%%%%
Let the following output be given as,
\begin{equation}
    y=\begin{bmatrix}
        \sin(x+\alpha)\\ \sin(x+\beta)
    \end{bmatrix}+D_w w
\end{equation}
The projection $\Pi$ to eliminate $D_w$ when applied to the output gives,
\begin{equation}
    \Pi y=\Pi\begin{bmatrix}
        \sin(x+\alpha)\\ \sin(x+\beta)
    \end{bmatrix}=\begin{bmatrix}
        \Pi_{11}\sin(x+\alpha)+\Pi_{12}\sin(x+\beta)\\
        \Pi_{12}\sin(x+\alpha)+\Pi_{22}\sin(x+\beta)
    \end{bmatrix}
\end{equation}
Without loss of generality choosing the first projected output via,
\begin{equation}
    \begin{bmatrix}I& 0\end{bmatrix}\Pi y=\Pi_{11}\sin(x+\alpha)+\Pi_{12}\sin(x+\beta)
\end{equation}
which is simplified into,
\begin{gather}
   \Pi_{11}\sin(x+\alpha)+\Pi_{12}\sin(x+\beta)=\\
    \left(\Pi_{11}\cos(\alpha)+\Pi_{12}\cos(\beta)\right)\sin(x)\\
    +\left(\Pi_{11}\sin(\alpha)+\Pi_{12}\sin(\beta)\right)\cos(x)
\end{gather}
and into,
\begin{equation}
   \begin{bmatrix}I& 0\end{bmatrix}\Pi y=A\sin(x+\theta)
\end{equation}
where 
\begin{equation}
    A=\sqrt{\left(\Pi_{11}\cos(\alpha)+\Pi_{12}\cos(\beta)\right)^2+\left(\Pi_{11}\sin(\alpha)+\Pi_{12}\sin(\beta)\right)^2}
\end{equation}
and
\begin{equation}
    \theta=\tan^{-1}\left(
        \frac{\Pi_{11}\sin(\alpha)+\Pi_{12}\sin(\beta)}{\Pi_{11}\cos(\alpha)+\Pi_{12}\cos(\beta)}\right)
\end{equation}
The following is stated,
\begin{equation}
    A \sin(x+\theta)=\begin{bmatrix}I& 0\end{bmatrix}\Pi y
\end{equation}
another solution is,
\begin{equation}
    A \sin(\pi-x-\theta)=\begin{bmatrix}I& 0\end{bmatrix}\Pi y
\end{equation}
Hence,
\begin{gather}
     x=-\theta+\sin^{-1}\left(\frac{1}{A}\begin{bmatrix}I& 0\end{bmatrix}\Pi y \right) \\
     x=-\theta+\pi-\sin^{-1}\left(\frac{1}{A}\begin{bmatrix}I& 0\end{bmatrix}\Pi y \right) 
\end{gather}



